% ============================================================================
%  GONDOLA INTELIGENTE - Mapeo de flujos del software
%  Compilar:  pdflatex flujos.tex
% ============================================================================
\documentclass[11pt,a4paper]{article}
\usepackage[utf8]{inputenc}
\usepackage[T1]{fontenc}
\usepackage[spanish]{babel}
\usepackage[margin=2cm]{geometry}
\usepackage{xcolor}
\usepackage{booktabs}
\usepackage{tikz}
\usetikzlibrary{shapes.geometric, arrows.meta, positioning, calc, fit}

\definecolor{capaf}{HTML}{E8F0FE}
\definecolor{proc}{HTML}{FFF8E1}
\definecolor{dec}{HTML}{FDECEC}
\definecolor{arch}{HTML}{E8F5E9}
\definecolor{bordec}{HTML}{37474F}

\tikzset{
  base/.style   = {font=\small, rounded corners=2pt, align=center},
  capa/.style   = {base, draw=bordec, very thick, fill=capaf},
  proc/.style   = {base, draw=bordec, thick, fill=proc},
  term/.style   = {base, draw=bordec, very thick, fill=dec,
                   rounded corners=6pt},
  datos/.style  = {base, draw=bordec, thick, fill=arch,
                   shape=cylinder, shape border rotate=90, aspect=0.25,
                   inner sep=3pt},
  dec/.style    = {base, draw=bordec, thick, fill=dec,
                   shape=diamond, aspect=1.6, inner sep=1pt},
  z/.style      = {base, draw=bordec, thick, fill=white},
  arr/.style    = {-{Stealth[length=3mm]}, thick, draw=bordec},
}

\newcommand{\etiqueta}[2]{\node[z] (#1) {#2};}

\begin{document}

\begin{center}
  {\LARGE\bfseries Góndola Inteligente}\\[2mm]
  {\large Mapeo de flujos del software (YOLO + FastAPI + React)}
\end{center}
\vspace{2mm}
\rule{\textwidth}{0.5pt}

\section{Arquitectura por capas}
Cada capa depende solo de las inferiores. Las flechas indican llamadas
``hacia abajo'' (la capa superior usa la inferior).

\begin{tikzpicture}[node distance=10mm]
  \node[capa] (intf) {
    \textbf{INTERFACES} \emph{(adaptadores)}\\
    \small cli.py (argparse) \quad api.py (FastAPI, endpoint REST)\\
    validar\_embudo.py (CLI de validación)\quad generar\_gt.py
  };
  \node[capa, below=of intf] (uc) {
    \textbf{USE CASES} \emph{(u orchestadores)}\\
    \small procesar\_video.py \quad exportar\_metricas.py
  };
  \node[capa, below=of uc] (d) {
    \textbf{DOMAIN} \emph{(reglas de negocio)}\\
    \small persona.py (máquina de estados) \quad analisis.py (zonas/heatmap)\\
    \small geometria.py (punto-en-polígono) \quad agente.py (recomendaciones)
  };
  \node[capa, below=of d] (inf) {
    \textbf{INFRASTRUCTURE} \emph{(drivers)}\\
    \small vision.py (YOLO pose + objetos, captura video)\\
    anonimizacion.py (blur facial)
  };
  \node[capa, below=of inf] (apl) {
    \textbf{APPLICATION / CONFIG}\\
    \small paths.py (rutas) \quad config.py (parámetros, zonas, umbrales)
  };
  \node[datos, right=38mm of uc] (json) {outputs/\\metricas\_retail.json};
  \node[datos, right=6mm of intf,xshift=10mm] (frames) {outputs/\\vision\_frames/\\f\_N.jpg};
  \foreach \a/\b in {intf/uc, uc/d, d/inf, inf/apl}{
    \draw[arr] (\a) -- (\b);
  }
  \draw[arr] (intf) -- node[above,font=\tiny] {lee/escribe} (json);
  \draw[arr] (uc) -- (json);
  \draw[arr] (uc) -- node[above,font=\tiny] {JPEG+índice} (frames);
\end{tikzpicture}

\section{Flujo principal del pipeline}
\emph{Entrada:} video (por defecto \texttt{video/videoBase.mp4} o subido)
$\rightarrow$ \emph{Salida:} \texttt{metricas\_retail.json} y fotogramas
anotados.

\begin{tikzpicture}[node distance=8mm]
  \node[term] (cli) {CLI: \texttt{python -m app.interfaces.cli}\\
    \tiny flags: --video, --no-ventana, --guardar-frames};
  \node[proc, below=of cli] (init) {DetectorPoseYObjetos()\\
    \tiny carga \texttt{yolo11n-pose.pt} + \texttt{yolov8n-producto-mano.pt}\\
    CapturaVideo() \tiny (lee fps / resolución)};
  \node[dec, below=of init] (loop) {¿quedan frames?};
  \node[proc, right=0mm of loop, xshift=40mm] (proc1) {si \texttt{PROCESAR\_CADA\_N}\\
    det.detectar\_objetos(frame)\\
    det.rastrear\_personas(frame)};
  \node[proc, below=of proc1] (proc2) {procesar\_personas()\\
    \tiny estados por id, eventos, zonas};
  \node[proc, below=of proc2] (render) {render: anonimizar, plot YOLO,\\
    polígonos de zona, etiquetas, heatmap,\\
    \tiny imwrite JPEG si --guardar-frames};
  \draw[arr] (cli) -- (init);
  \draw[arr] (init) -- (loop);
  \draw[arr] (loop) -- node[above] {sí} (proc1);
  \draw[arr] (proc1) -- (proc2);
  \draw[arr] (proc2) -- (render);
  \draw[arr] (render) -| ($(render.south -| loop.south) + (0,-14mm)$);
  \draw[arr,dashed] ($(render.south -| loop.south) + (0,-14mm)$) -- (loop.south);
  \node[proc, below=of loop, yshift=-14mm] (fin) {fin: exportar\_json\_metricas()\\
    \tiny dedupe embudo, zonas calientes/frías, heatmap, recomendaciones};
  \node[datos, right=30mm of fin] (cli2) {metricas\_retail.json};
  \draw[arr] (loop) -- node[above] {no} (fin);
  \draw[arr] (fin) -- (cli2);
  \node[proc, left=of fin, xshift=-12mm] (idx) {escribe vision\_frames\_index.json\\
    \tiny \{fps, stride, frames\}};
  \draw[arr] (fin.west) -- (idx.east);
\end{tikzpicture}

\section{Máquina de estados de la persona (\texttt{clasificar})}
Estados: \texttt{En Tránsito}, \texttt{Alcanzar Estante}, \texttt{Inspeccionar
Producto}, \texttt{Mano en Estante}, \texttt{Devolver el producto},
\texttt{Producto Adquirido}.

\begin{tikzpicture}[node distance=12mm]
  \node[proc] (z1) {En Tránsito};
  \node[proc, above=of z1, xshift=-16mm] (z2) {Alcanzar Estante};
  \node[proc, above=of z2, xshift=-16mm] (z3) {Inspeccionar Producto};
  \node[proc, right=of z2] (z4) {Mano en Estante};
  \node[proc, below=of z1, xshift=-16mm] (z5) {Devolver el producto};
  \node[proc, right=of z1] (z6) {Producto Adquirido};
  \draw[arr] (z1.90) to[out=90,in=270] node[left,pos=.25] {\tiny entra a zona} (z2.270);
  \draw[arr] (z2.90) to[out=90,in=270] node[left,pos=.25] {\tiny quieto $>$ \texttt{QUIETO\_PARA\_VER}} (z3.270);
  \draw[arr] (z3.east) to[out=0,in=90] node[above,pos=.4] {\tiny muñeca en agarro} (z4.90);
  \draw[arr] (z2.east) to[out=0,in=180] node[below,pos=.25] {\tiny toma (pick-up) o retirado} (z6.180);
  \draw[arr] (z1.east) to[out=0,in=180] node[above,pos=.4] {\tiny sale de zona} (z4.270);
  \draw[arr] (z4.east) to[out=0,in=90] node[right,pos=.4] {\tiny sale de zona / sin objeto} (z6.0);
  \draw[arr] (z5.east) to[out=0,in=270] node[below,pos=.4] {\tiny devolución} (z6.270);
  \draw[arr] (z2.south) to[out=270,in=90] node[right,pos=.3] {\tiny retirado $\wedge$ sin objeto} (z5.90);
  \draw[arr] (z3.east) to[out=0,in=90] node[above,pos=.4] {\tiny item comprado/quieto} ([xshift=10mm]z4.east);
\end{tikzpicture}

\subsection{Decisiones clave de \texttt{clasificar} (seudocódigo)}
\begin{itemize}
  \item \textbf{Fuera de zona}: si \texttt{adquirido} $\rightarrow$ sigue
  Adquirido. Si \texttt{retirado} (o mano en estante hace poco +
  \texttt{tiene\_objeto}) $\rightarrow$ \textbf{pick-up}, asociado a la
  \texttt{ultima\_zona}. Si no, En Tránsito.
  \item \textbf{Dentro de zona}: si \texttt{adquirido} $\rightarrow$
  Adquirido. Si Inspeccionando y quieto $\rightarrow$ \texttt{compro=True} y
  pick-up inmediato. Si \texttt{devolviendo} $\rightarrow$ Devolver. Si
  \texttt{mano\_dentro} $\rightarrow$ Mano en Estante. Si lleva tiempo quieto
  $\rightarrow$ Inspeccionar Producto. Si no, Alcanzar Estante.
\end{itemize}

\section{Detección de interacción mano--producto (\texttt{detectar\_interaccion\_manos})}

\begin{tikzpicture}[node distance=8mm]
  \node[dec] (g1) {¿zona=None o\\adquirido?};
  \node[proc, below=of g1] (p1) {posesion = estricta o relajada\\
    \tiny (muñeca $\leftrightarrow$ caja de objeto)};
  \node[dec, below=of p1] (g2) {¿flanco posesión\\
    False$\rightarrow$True?};
  \node[proc, below=of g2] (p2) {\texttt{retirado=True}; \texttt{alcance=True}\\
    \tiny pick-up preparado (POSESION\_DISPARA\_PICKUP)};
  \node[dec, below=of p2, yshift=-6mm] (g3) {¿muñeca confiable\\en polígono \texttt{agarro}?};
  \node[proc, below=of g3] (p3) {mano\_en\_zona++ ;\\\tiny mano\_estante\_frame = fi};
  \node[dec, below=of p3] (g4) {¿mano\_en\_zona $\ge$ 5\\y no mano\_dentro?};
  \node[dec, below=of g4] (g5) {¿retirado y\\sin objeto?};
  \node[proc, below=of g5] (p5) {\texttt{devolviendo=True} (put-back)};
  \node[proc, right=22mm of g5] (p6) {\texttt{mano\_dentro=True}\\
    \tiny event ``Mano en Estante''};
  \node[proc, right=22mm of g4] (p7) {flanco objeto+mano\_dentro\\
    $\rightarrow$ \texttt{retirado=True}};
  \node[proc, left=22mm of g3, xshift=-6mm] (p8) {si mano\_dentro y objeto\\
    \tiny $\rightarrow$ retirado; reset mano\_en\_zona};
  \draw[arr] (g1) -- node[right] {\tiny no} (p1);
  \draw[arr] (p1) -- (g2);
  \draw[arr] (g2) -- node[right] {\tiny sí} (p2);
  \draw[arr] (g2.south) -- (g3.north);
  \draw[arr, dashed] (p2.south) to[out=270,in=90] (g3.east);
  \draw[arr] (g3) -- node[right] {\tiny sí, en agarro} (p3);
  \draw[arr] (g3) -- node[above] {\tiny fuera} (p8);
  \draw[arr] (p3) -- (g4);
  \draw[arr] (g4.west) -- node[above] {\tiny aún no} (p7);
  \draw[arr] (g4.south) -- node[right] {\tiny sí} (g5);
  \draw[arr] (g5.south) -- node[right] {\tiny sí} (p5);
  \draw[arr] (g5.west) -- node[above] {\tiny no (tiene objeto)} (p6);
  \node[term, below=of p5, yshift=2mm] (inc) {fin del gesto:\\
    \tiny clasificar() convierte retirado$\rightarrow$pick-up\\
    o devolviendo$\rightarrow$put-back};
  \draw[arr] (p5.south) -- (inc.north);
  \draw[arr] (p6.south) -| (inc.east);
\end{tikzpicture}

\subsection{Señales de posesión (\texttt{detectar\_posesion\_objeto} / \texttt{detectar\_posesion\_relajada})}
\begin{itemize}
  \item \textbf{Estricta}: una muñeca con confianza $\ge$ \texttt{CONF\_MANO}
  solapa o queda a $\le$ \texttt{MANO\_OBJ\_TOLERANCIA} (escalado a la
  resolución real) de una caja de producto.
  \item \textbf{Relajada}: distancia $\le$ \texttt{MANO\_OBJ\_TOLERANCIA\_RELAJADA}
  (\texttt{= tolerancia + radio}). Captura el agarre en planos frontales donde
  los dedos no cubren la caja.
  \item La zona del gesto se decide \textbf{priorizando la muñeca} (agarro)
  sobre el centroide (polígono).
\end{itemize}

\section{Exportación de métricas (\texttt{exportar\_json\_metricas})}
\begin{tikzpicture}[node distance=8mm]
  \node[term] (e1) {fin del video};
  \node[proc, below=of e1] (e2) {dedupe global pick-ups / put-backs\\
    \tiny misma zona en \texttt{VENTANA\_DEDUPE\_PICKUP} frames};
  \node[proc, below=of e2] (e3) {por persona: eventos + trayectoria\\
    \tiny segundo (frame/fps), x/y normalizados};
  \node[proc, below=of e3] (e4) {analizador.resumen(fps): embudo,\\
    clasificación caliente/templada/fría,\\
    visitas, dwell, presencia, rechazo};
  \node[proc, below=of e4] (e5) {heatmap.exportar\_grilla(celdas)\\
    \tiny grilla global para dashboard};
  \node[proc, below=of e5] (e6) {agente determinista:\\
    \tiny recomendar\_zona + recomendar\_planograma};
  \node[datos, right=34mm of e3] (j) {metricas\_retail.json\\
    \tiny embudo\_interaccion, analisis\_zonas,\\
    detalle\_por\_persona, recomendaciones};
  \draw[arr] (e1) -- (e2) -- (e3) -- (e4) -- (e5) -- (e6);
  \draw[arr] (e3) -- (j);
  \draw[arr] (e4) -- (j);
  \draw[arr] (e5) -- (j);
  \draw[arr] (e6) -- (j);
\end{tikzpicture}

\section{Validación frente a ground truth}
\begin{tikzpicture}[node distance=8mm]
  \node[proc] (v1) {\texttt{generar\_gt.py}\\ \tiny vuelca candidatos.json desde los eventos detectados};
  \node[proc, below=of v1] (v2) {\textbf{editor}: revisar/ajustar\\
    \tiny \texttt{ground\_truth/candidatos.json}};
  \node[proc, below=of v2] (v3) {\texttt{validar\_embudo.py --gt candidatos.json --barrido-tol}\\
    \tiny match vecino más cercano (¡criterio greedy, listas ordenadas!)};
  \node[proc, below=of v3] (v4) {métricas TP/FP/FN\\(tol $\pm$15 por defecto), F1 por zona,\\barrido de tolerancias (2..31), umbral F1};
  \node[term, below=of v4] (v5) {F1 y sanity $\Rightarrow$ decide el umbral\\
    \tiny ¿aprueba el modelo? (--umbral-f1)};
  \draw[arr] (v1) -- (v2) -- (v3) -- (v4) -- (v5);
\end{tikzpicture}

\section{API y dashboard (secuencia)}
\begin{tikzpicture}[node distance=8mm]
  \node[proc] (fe1) {Usuario: Subir video};
  \node[proc, below=of fe1] (fe2) {POST /api/video\\
    \tiny valida extensión + tipo MIME, guarda en video/subidos};
  \node[proc, below=of fe2] (fe3) {spawn: \texttt{cli --no-ventana --video X --guardar-frames outputs/vision\_frames}\\
    \tiny thread background + lock anti-concurrencia};
  \node[datos, below=of fe3] (fe4) {metricas\_retail.json\\ + vision\_frames/*.jpg + índice};
  \node[proc, below=of fe4] (fe5) {GET /api/datos \quad GET /vision-frames/*\\
    \tiny normaliza zonas (ids crudos), 404 del backend $\Rightarrow$ /datos.json demo};
  \node[proc, below=of fe5] (fe6) {React App.jsx: mergeZones() + visión de góndolas\\
    \tiny intervalo = stride×1000/fps, fotograma i/120};
  \draw[arr] (fe1) -- (fe2) -- (fe3) -- (fe4) -- (fe5) -- (fe6);
\end{tikzpicture}

\section{Endpoints de la API}
\begin{center}
\begin{tabular}{@{}lll@{}}
\toprule
\textbf{Método} & \textbf{Ruta} & \textbf{Efecto}\\
\midrule
GET & /api/estado & estado de las métricas y análisis en curso\\
GET & /api/datos & métricas normalizadas (o fallback demo)\\
POST & /api/analizar & dispara el pipeline en background (opcional ruta)\\
POST & /api/video & sube video $\rightarrow$ guarda $\rightarrow$ analiza\\
GET & /api/vision-frames & índice JSON \{fps, stride, frames[]\}\\
GET & /api/vision-frames/* & JPEGs anotados (secuencia del visor)\\
GET & / & frontend estático (dist/ o public/)\\
\bottomrule
\end{tabular}
\end{center}

\section{Constantes relevantes (app/config.py)}
\begin{center}
\begin{tabular}{@{}ll@{}}
\toprule
\textbf{Constante} & \textbf{Efecto}\\
\midrule
\texttt{PROCESAR\_CADA\_N=2} & stride de procesamiento (frames saltados)\\
\texttt{CONF\_OBJETOS} & umbral de confianza del detector de producto en mano\\
\texttt{MANO\_RADIO} / \texttt{MANO\_OBJ\_TOLERANCIA} & posesión estricta (solape/distancia muñeca--caja)\\
\texttt{MANO\_OBJ\_TOLERANCIA\_RELAJADA} & radio de alcance (tolerancia + radio)\\
\texttt{POSESION\_DISPARA\_PICKUP} & flanco de posesión dispara la toma sin `agarro`\\
\texttt{MANO\_EN\_ZONA\_PARA\_RETIRAR=5} & frames mínimos de mano en agarro\\
\texttt{VENTANA\_DEDUPE\_PICKUP=25} & dedupe de eventos por zona\\
\texttt{VENTANA\_SALIDA\_OBJETO=8} & pick-up al salir de la zona con objeto reciente\\
\texttt{IMGSZ=960} / \texttt{CONF=0.25} & pose\\
\bottomrule
\end{tabular}
\end{center}

\end{document}